\chapter{Integral Calculus }
  \label{chapt:integration}
  \markboth{{\sc Integration}}{{\sc Integration}}

\section{Introduction}  
  The first section of this book dealt with dealt with \emph{Differential
  Calculus}, literally the rules for differences.  We learned how to
  differentiate various functions and how to apply this to various
  modeling problems.  For example, one of the exercises we had was to
  show that a hanging chain (catenary) satisfied the differential
  equation
  $$
  \dfdxn{y}{x}{2} = \frac{W}{H}\sqrt{1+\left(\dfdx{y}{x}\right)^2}.
  $$
  A later exercise had you show that
  $$
  y=\frac{H}{w}\left(\frac{e^{\frac{wx}{H}}+e^{-\frac{wx}{H}}}{2}\right).
  $$
In fact, satisfied this differential equation.  All this is well and
good, but you might have been wondering at the time where we pulled
this function out of to differentiate.  This is what integral calculus
is about.  In general, given a function $y=y(x),$ we can differentiate
it to obtain the differential $\d{y}.$  Integration is the process of
starting with a differential and finding the function that was
differentiated to obtain it.

For example, if $y=x^2$ then $\d{y}=2x \d{x}.$ It stands to reason
that if we integrate $\d{y}=2x\d{x},$ we should end up with $y=x^2.$
There is a little bit of a complication here, as differentiating
$y=x^2+c$ yields $\d{y}=2x$ dx for any constant $c.$ It turns out that
this is the only such complication; if $\d{f}=\d{g},$ then $f$ and $g$
must differ by a constant.  To see why, notice that if $\d{f}=\d{g}$
then $\d{(f-g)}=\d{f}-\d{g}=0.$ Dividing by $\d{x}$ so we can think of
this as a slope, we have $\dfdx{(f-g)}{x}=0.$ What kind of curve would
have a slope which is constantly $0?$ A moment's thought will suggest
that this curve must be a horizontal line Indeed, if it was not a
horizontal line, then somewhere it should have a nonzero slope.  Such
an argument can be made rigorous, but for now, we will go with this
intuitive understanding.

With this place, we will introduce some notation.  The notation is
motivated by the fact that integration (putting together, summing) is
the inverse process of differentiation (taking a difference).  With
this in mind, the integral of a differential $dy$ is denoted by  
$$
\int\d{y}
$$
where the ``stylized S'' really means a sum of differentials
(differences).  There are a couple of things to note about this
notation.  First, \underline{\bf{}we only integrate differentials.}
In other words, whatever you are integrating must be a differential.
For example, $\int 2x \d{x}$ makes sense, $\int 2x$ does not.

Second, notice we did not write $\int\d{y}=y.$  This is because $\int \d{y}$ has an
infinite number of possibilities.  Fortunately, as we mentioned
before, these all differ by constants.   Thus we can write something
like 
$$
\int 2x \d{x}=x^2+c
$$
where $c$ is an arbitrary constant.  This represents all of the
possible functions whose derivative is $2x,$ so with that in mind
$\int 2x \d{x}$ is referred to as the indefinite integral or
anti-derivative of $2x.$ These names come from the foundational switch
from differentials to limits and derivatives, but we still want to
emphasize that we integrate differentials, not functions (so the
$\d{x}$ is not optional).

\section{Rules for Integration}
\label{sec:rules-integration}
Just as Differential Calculus is literally the rules for (infinitely
small) differences, Integral Calculus is the rules for summing up
these differences.  Not surprisingly, each differentiation rule has
its integration rule counterpart.  The following table illustrates
some of these rules.
\begin{wraptable}{r}{5in}
  \begin{tabular}{ |c|c| }
    \hline
    \multicolumn{2}{|c|}{\bf{}General} \\
    \multicolumn{2}{|c|}{\bf{}Integration Rules} \\\hline
    Differentiation Rule & Integration Rule\\\hline
    $\d{(\text{constant})}=0$&$\int 0\d{x}=C,$ where $C$ is a
                               constant\\\hline
    $\d(f+g)=\d{f}+\d{g}$&
                           $\int\d(f+g)=\int\d{f}+\d{g}=f+g+c=\int\d{f}+\int\d{g}$\\\hline
    If a is a constant, then $\d(ay)=a \d{y}$& $\int d(ay)=\int a\d{y}=ay+c=a\int\d{y}$\\\hline{}
$\d{x}^{n+1} =(n+1) x^n \d{x}$ & $\int x^n
                                 \d{x}=\frac{x^{n+1}}{n+1}+c, \text{
                                 where }n\ne-1$\\\hline
    $d(\ln(x))=\frac{1}{x} \d{x}=x^{-1} \d{x}$&	$\int x^{-1}
                                               \d{x}=\ln\abs{x}+c$\\\hline{}
    $\d(e^x)=e^x\d{x}$&	$\int e^x \d{x})=e^x+c$\\\hline{}
    $\d(\sin(x))=\cos(x)  \d{x} $&$\int(\cos(x) \d{x}=\sin(x)+c$\\\hline{}
    $\d(\cos(x) )=-\sin(x)\d{x}{}$& $\int \sin(x) \d{x}=-\cos(x)+C$\\\hline
    $\d(\tan(x))=\sec^2(x)  \d{x}$& $\int\sec^2(x) \d{x}=\tan(x)+C$\\\hline
    $\d(\sec(x))=\sec(x) \tan(x) \d{x}$& $\int \sec(x) \tan(x) \d{x}=\sec(x)+C$\\\hline
    $\d(\inverse{\sin}(x) )=\frac{1}{\sqrt{1-x^2}}  \d{x}$& $\int\frac{1}{\sqrt{1-x^2}}\d{x}=\inverse{\sin}(x)+c$\\\hline
    $\d(\inverse{\tan}(x))=\frac{1}{1+x^2}  \d{x}$&$\int\frac{1}{1+x^2}  \d{x}=\inverse{\tan}(x)+C$\\\hline
    $\d(\inverse{\sec}(x))=\frac{1}{x\sqrt{x^2-1}}  \d{x}$&$\int \frac{1}{x\sqrt{x^2-1}}  \d{x}=\inverse{\sec}(x)+C$\\\hline{}
  \end{tabular}
\end{wraptable}

This table is not exhaustive, as every differentiation performed gives
rise to a corresponding integration formula.  Conspicuously absent are
the product rule, quotient rule, and chain rule.  The analog of the
quotient rule is rather limited in its applicability and typically is
not used.  The analogs of the others are quite useful and will be
covered later.

Notice that the analog of the power ruler (line 4)
makes sense verbally.  To differentiate to a power, you multiply by
the exponent and then subtract one from the exponent.  To reverse the
process, you reverse each step and reverse the order of the steps
(much as you put your socks on and then your shoes and to reverse it
you take off your shoes first and then your socks).  Thus you add one
to the exponent and then divide by the exponent.  Notice that the rule
will not work if the exponent is zero as you cannot divide by zero.  This
is where the natural logarithm fills in the gap (line 5).  While we
are at line five, notice that we have inserted an absolute value:
$\int\frac1x\d{x}=\ln\abs{x}+C.$

We hand't brought it up before, but the logarithm of a negative number
is not a real number.  It turns out that inserting an absolute value
allows us to extend the natural logarithm function to negative numbers
while still maintaining the same differentiation rule.

\begin{embeddedproblem}{}
  One way to define the absolute value function is as follows
  $$
  \abs{x}=\sqrt{x^2}.
  $$
  where this signifies the non-negative square root.  For example,
  $\abs{-3}=\sqrt{(-3)^2}=\sqrt{9}=3$ and
  $\abs{3}=\sqrt{(3)^2}=\sqrt{9}=3$.  Show that with this definition
  of absolute value, we still obtain the formula
  $\d(\ln\abs{x})=\frac{1}{x}\d{x}.$

Note:  You still must be careful that you do not try to write the logarithm of a negative number.
\end{embeddedproblem}

A common practice is to write $\int \inverse{x} \d{x}=\ln(x)+C,$ , even
though this is only valid for .  We will adopt this practice, and
resort to an absolute value only when there is an issue about being
negative.

\begin{embeddedproblem}{}
  It looks like we broke our own rule in lines two and three in the above table of integration rules.  After all,
$$
\int \d{f}+int \d{g}=f+c_1+g+c_2
$$
where $c_1$ and $c_2$ are arbitrary constants, and 
$$
a \int \d{y}=a(y+c)=ay+ac.
$$
Explain why what we have in the table is still legitimate.
\end{embeddedproblem}

In light of the above problem, practitioners often will perform the
integration of a differential applying whatever tools are appropriate
and then combining any arbitrary constants into one constant at the
end.  For example, consider 
\begin{align*}
  \int \left((5x^2-3\sqrt{x}+\frac{2}{1+x^2}\right)  \d{x} 
  &=5\int x^2\d{x}-3\int x^{1/2}  \d{x}+2\int\frac{1}{1+x^2}  \d{x}\\
    1&=5\left(\frac{x^3}{3}\right)-\frac{3x^{3/2}}{\frac{3}{2}}+2 \inverse{\tan}(x)+c\\
  &=\frac{5}{3} x^3-2x^{\frac{3}{2}}+2 \inverse{\tan}(x)+C
\end{align*}

As with differentiation rules, when you get more comfortable with the
integration rules, you will probably do these in your head.  Again,
developing this fluency this is why you need to practice. 

\begin{PracticeProblemSection}
Integrate the following
\end{PracticeProblemSection}

We have made a good start, but we are still a bit away from solving differential equations such as 
  $$
  \dfdxn{y}{x}{2} = \frac{W}{H}\sqrt{1+\left(\dfdx{y}{x}\right)^2}.
  $$

A start might be to let  z=dy/dx, so we have
\begin{align*}
\dfdx{z}{x}&=\frac{w}{H} \sqrt{1+z^2}\\
\frac{1}{\sqrt{1+z^2}}  \d{z}&=\frac{w}{H}  \d{x}\\
\int\frac{1}{\sqrt{1+z^2}}  \d{z}&=\int\frac{w}{H}  \d{x}
\end{align*}

The right-hand side of this doesn't pose any problems since $\frac{w}{H}$ is a
constant.  However, the left-hand side is another matter, since it
doesn't quite fit with anything in our table.  We will learn
techniques to handle such situations, but first let's look at some
applications of integration which utilize what we know. 

\section{\sc{}First application --  projectile motion}

Suppose you throw a ball straight up into the air.  Of course, it will
go up and come back down.  For the moment, let's ignore air
resistance.  How high will the ball go?  How long will it take the
ball to hit the ground?  What will be its impact velocity?  Suppose
you throw the ball up with twice the initial speed.  Will the ball go
twice as high?  These were the types of questions that interested
scientists like Galileo Galilei $(1564-1642)$ who were some of the first
scientists to systematically apply mathematics to science,
specifically projectile motion. 

\textcolor{red}{Bud -- I think it makes sense to put the stuff on motion here --
specifically, Galileo deducing that the acceleration due to gravity is
constant -- Section $5.1.$  I will continue based off of that. }

  
Now that we have that acceleration due to gravity is constant, let's
look more into projectile motion.  This fits into our discussion on
integral calculus since we are given information about the
acceleration and must integrate twice to recover the position of the
ball.  Specifically, let's say that the ball is thrown upward with an
initial velocity of $19.6$ meters per second (roughly $44$ miles per hour)
and is released at a point $2$ meters above the ground.  How high will
the ball go?  When will it hit the ground?  What will be the impact
velocity?

To answer these questions, we must translate this problem into a mathematics problem.  
With this in mind, suppose the height of the ball is given by $y=y(t)$ where $y$ is measured in meters and $t$ is measured in seconds.  Then we have the following information
$$
\dfdxn{y}{x}{2}=-9.8 \frac{\text{m}}{\text{s}^2},\hfill
\left.\dfdx{y}{t}\right|_{t=0}=19.6 \frac{\text{m}}{\text{s}},\hfill
y(0)=2 \text{m}.
$$

Notice that the acceleration is constant (we are only looking a
gravity and ignoring air resistance) and is negative (we've set up $y$ 
so that up is the positive direction and down is the negative
direction).

With this in mind, we want to find the following quantities:
$y$ when  $\dfdx{y}{t}=0$ (how high),    $t$ when $y=0$  (when it hits the
ground), $\left. \dfdx{y}{t}\right|_{y=0}$   (impact velocity).

For convenience, let $v=\frac{y}{t}$ denote velocity, and
$a=\dfdx{v}{t}=\dfdxn{y}{t}{2}$ denote acceleration.  Starting with our constant
acceleration $a=-9.8,$ we need to integrate to find $v$ and $y.$  Once we
find these, then we can answer the questions posed. 
\begin{align*}
a&=\dfdx{v}{t}=-9.8\\
\d{v}&=-9.8 \d{t}\\
\int \d{v}&=\int-9.8 \d{t}\\
v+c_1&=-9.8t+c_2\\
v&=-9.8t+(c_2-c_1)
\end{align*}

Since $c_1$ and $c_2$ are arbitrary constants, then we will rewrite
$c=c_2-c_1$ and notice that $c$ is an arbitrary constant.  From this point
on, we will automatically combine arbitrary constants whenever we can
without saying so.  To find what $c$ is in this case, we utilize the
fact that $v(0)=19.6.$  Thus 
\begin{align*}
19.6&=v(0)=-9.8(0)+c\\
19.6&=c\\
\end{align*}
and so our velocity is given by 
$$
v=-9.8t+19.6.
$$
We now integrate this to find the ball's position.
\begin{align*}
\dfdx{y}{t}&=-9.8t+19.6\\
\d{y}&=(-9.8t+19.6)\d{t}\\
\int\d{y}&=\int(-9.8t+19.6)\d{t}\int
y&=-9.8\frac{t^2}{2}+19.6t+C.
\end{align*}

To find the constant $c$ now, we utilize that $y(0)=2.$ 
\begin{align*}
  2&=y(0)=-4.9(0)^2+19.6(0)+C\\
  2&=c
\end{align*}

Thus the ball's positon is given by $y=-4.9t^2+19.6t+2.$ 
Once we have these three equations of motion (acceleration, velocity,
position) we can answer any question about the ball's motion.  In
general, we suggest you find these three equations before trying to
answer any questions about motion.

To find the maximum height, we set $v=0.$ 
\begin{align*}
0&=-9.8t+19.6\\
t&=2 \text{sec}\\
y(2&)=-4.9(2^2 )+19.6(2)+2=21.6 \text{meters}
\end{align*}
so the ball will reach a height of $21.6$ meters.

To find the impact velocity, we first need to find when the ball hits
the ground, or $t$ when $y=0.$ 
\begin{align*}
0&=y=-4.9t^2+9.6t+2\\
t&=\frac{-9.6\pm\sqrt{9.6^2-4(-4.9)2}}{2(-4.9)} \\
\end{align*}
so $t=4.1$ or $t=-0.1$ 
Since the negative time would represent time in the past, we can say
that the ball will hit the ground when $t=4.1$ sec.  Thus the impact
velocity is  
$$
v(4.1)=-9.8(4.1)+19.6=-20.58 \text{m/s.}
$$
The negative answer makes sense as the ball is traveling in the
negative direction (given how we set up the $y$ axis). 


\begin{embeddedproblem}{}
  Suppose we throw a ball straight up with an
  initial velocity of $v_0$ and then throw it up with an initial
  velocity of $2v_0?$  Will the ball travel twice as high (from the
  release point)?  Again, we will ignore air resistance.
\end{embeddedproblem}

\begin{PracticeProblemSection}

\end{PracticeProblemSection}

Of course, the same principles will work in two dimensions as with a
single dimension.  All we do is treat the horizontal and vertical
directions separately.  With this in mind, the horizontal acceleration
would be zero. 
\begin{embeddedproblem}{}
  Consider the following projectile launched from the origin, with an
  initial velocity of $v_0$ and angle of elevation $\theta.$ \\
\includegraphics*[height=2in,width=5in]{../Figures/PIC-Projectile}.

\begin{enumerate}[label={\bf{}(\alph*)}]
\item 	Show that the trajectory is given by the parabola
  $y=-\frac{g}{2}{v_0  cos(\theta)^2} x^2+\tan(\theta)x.$
\item 	Use the equation in question $1$ to show that the range of the
  projectile is given by  $\frac{v_0^2}{g}  \sin(2\theta) . $
\item 	Show that the maximum range occurs when
  $\theta=\frac{\pi}{4}.$ 
\item 4.	On Feb. 6, 1971, astronaut Alan Shepard pulled out a
  makeshift six-iron he smuggled on board Apollo $14$ and hit two golf
  balls on the lunar surface, becoming the first, and only, person to
  play golf anywhere other than earth.  Assuming that the gravity on
  the moon is $1/6$ that of earth and that Shepard can hit the ball
  identically on the earth and the moon, and ignoring air resistance
  on the earth, would the ball go $6$ times farther before it touches
  the ground on the moon than on the earth? Justify your answer. 
\end{enumerate}
\end{embeddedproblem}

{\noindent{}\bf{}Including Air Resistance -- This Looks Like a Job for
  Substitution!}
Previously we examined the maximum height of a ball thrown straight up
with an initial velocity of $v_0$ where air resistance was ignored.  We
will now consider air resistance in the problem.

With this in mind, we will again let $y=y(t)$ denote the height of the
ball at time $t,$ $v=dy/dt$ the velocity of the ball, and
$a=\dfdx{v}{t}= \dfdxn{y}{t}{2}$ the acceleration of the ball.
      
As we said, we will now consider air resistance.  One of the tenants
of mathematical modeling is to start simple.  With that in mind, let's
assume that the force due to air resistance is proportional to the
velocity, which is reasonable at relatively low speeds and has the
advantage of being relatively simple. You can experience this for
yourself by moving your hand through the air.  As you move your hand
faster, you notice more air resistance.  For higher speeds where
turbulence is more of a factor, it is actually more accurate to model
air resistance as proportional to the square of the velocity, but we
will stick with this simpler linear model.  As such, we will assume
that the force due to air resistance is given by $-Rv,$ where $R>0$ is
a constant (which depends on the viscosity of the air, the
aerodynamics of the ball, etc.).  Note we have $-R$ since the force
due to air resistance is the opposite direction of the motion of the
ball.
        
We will also assume that the acceleration due to gravity is given by
$-g,$ where $g>0$ is constant, and the mass of the ball is $m.$ (We
are assuming our axis $y$ is labeled with the positive direction being
up.)

\begin{embeddedproblem}{}
  \begin{enumerate}[label={\bf{}(\alph*)}]
  \item Utilizing that $\text{force}=\text{mass}\text{acceleration,}$
    show that the velocity of the ball must satisfy the equations 
    \begin{align*}
      m \dfdx{v}{t}&=-mg-Rv\\
      v(0)&=v_0
    \end{align*}
and use this to show that
$$
\frac{\d{v}}{mg+Rv}=-\frac{1}{m} \d{t}.
$$
\item 	Use the substitution $z=mg+Rv,$ to rewrite the above equation as
$$
\frac{\frac{1}{R} \d{z}}{z}=-\frac{1}{m} \d{t}
$$
and show that 
$$
v=\left(\frac{mg}{R}+v_0\right) e^{-\frac{Rt}{m}}-\frac{mg}{R}.
$$
\item 	Assume that the initial position of the ball is given by $y(0)=0,$ and show that the height of the ball at time $t$ is given by
$$
y=\frac{m}{R} 
\left(
  \frac{mg}{R}+v_0 
\right)
\left(
  1-e^{-\frac{Rt}{m}}
\right) -\frac{mg}{R} t
$$
\item 	Show that the maximum height of the ball can be written as
$$
\frac{gm^2}{R^2}
\left(
  \frac{Rv_0}{gm}-ln\left(1+\frac{Rv_0}{gm}\right) 
\right).
$$
\item Show that the formula for the maximum height given in part (d)  is
  positive for any positive values of $g, m, R,$ or $ v_0.$ [Hint:  Find the
  tangent line to the curve $y=\ln(1+x)$ at $x=0.$  Graph these!  Which is
  bigger and what does this have to do with the formula in part (d)?] 
  \end{enumerate}
\end{embeddedproblem}

The above problem utilized the technique of substitution and is the
integration analog of the chain rule, which itself was the technique
of substitution applied to differentiation.  As with differentiation,
this did not actually compute the integral, but made it simpler to
apply formulas and techniques with which we are already familiar.
This technique of substitution is very powerful and allows us to
integrate differentials which are not in our table.  We already saw
one example in our previous problem.  For another, consider that in
the table, we had $\int \sin(x) \d{x}$ and $\int\cos(x)\d{x},$ but we did not
have $\int\tan(x)\d{x}.$  Now we have the tools to do this one.         
\begin{myexample}
  \label{ex:int-tan}
  If we make the substitution $z=\cos(x),$ then $\d{z}=-sin(x) \d{x}$
  and so $\d{x}=-\frac{\d{z}}{sin(x)}.$

  Thus
  \begin{align*}
    \int \tan(x) \d{x}&=\int \frac{\sin(x)}{\cos(x)} \d{x}\\
                      &=\int\frac{\sin(x)}{z}\left(-\frac{\d{z}}{\sin(x)}\right)\\
                      &=-\int  \frac{\d{z}}{z}\\
                      &=-\ln(z)+C\\
                      &=-\ln(\cos(x))+C\\
                      &=\ln\left(\frac{1}{\cos(x)} \right)+C\\
    &=\ln(\sec(x))+C      
  \end{align*}
\end{myexample}

\begin{embeddedproblem}{}
  Compute $\int\cot(x)\d{x}$ using the same technique as in
  Example~\ref{ex:int-tan}.
\end{embeddedproblem}

The beauty of a substitution is that there is no such thing as a wrong
substitution.  For example, in the above computation of $\int  \tan(x)
\d{x}$, we could have let $z=x,$ so that $\d{z}=\d{x},$ and we would
have $\int \tan(x)  \d{x}=\int \tan(z)  \d{z},$ which though correct,
would not have been very helpful.  Alternatively, we could have let
$u=tanx$ so that $\d{u}=\sec^2(x)  \d{x},$ and so 
$$
\int \tan(x)  \d{x}=\int \frac{u \d{u}}{\sec^2(x)} =\int
\frac{u}{1+tan^2x}  \d{u}=\int \frac{u}{1+u^2}  \d{u}.
$$
This is a doable integral, but this seems worse than the original
strategy.  Hey!  If we originally had the integral $\int \frac{u}{1+u^2} \d{u},$
we could utilize the same substitution, to transform this into the
integral $\int \tan(x) \d{x},$ which we now know how to do.  This would be
entirely correct, but alas, it is doing this integral the hard way.
We could compute $\int \frac{u}{1+u^2} \d{u},$ more easily by letting
$y=1+u^2,$ so that $\d{y}=2u \d{u},$ and so
$$
\int \frac{u}{1+u^2}  \d{u}=\int \frac{u}{y}  \left(\frac{\d{y}}{2u}\right)=\frac{1}{2} \int \frac{\d{y}}{y}=\frac{1}{2}\ln(y)+C=\ln\sqrt{1+u^2 }+C.
$$
The lesson in all of this is that the good news is that you can do any
substitution you wish and as long as it is performed correctly, it is
not wrong.  The bad news is that you can do any substitution you wish,
and as long as it is performed correctly, it is not wrong, but it may
not be useful.  Sometimes they can not only lead to harder integrals,
but ones where there is nothing you can really do (at least readily).

For example, consider
$$
\int x^2 \sqrt{x^3+1} \d{x}.
$$

We could let $u=x^2$ so that $\d{u}=2x$ dx.  This would lead to 
$$
\int x^2 \sqrt{x^3+1} \d{x}=\int u\sqrt{u^{3/2}+1}
\left(\frac{\d{u}}{2x}\right)=\frac{1}{2} \int
\frac{u\sqrt{u^{3/2}+1}}{x}  \d{u}.
$$

Notice that we can ``bring the $\frac{1}{2}$ outside the integral''
because it is a constant, but we cannot do the same with
$\frac{1}{x},$ because it is not a constant (it depends on $u$) and
that is not a legal move.  We could rewrite  $x=\sqrt{u},$ so that 
$$
\int x^2 \sqrt{x^3+1} \d{x}=\frac{1}{2} \int
\frac{u\sqrt{u^{3/2}+1}}{\sqrt{u}}  \d{u}=\frac{1}{2} \int
\sqrt{u}\sqrt{u^{3/2}+1} \d{u}=\frac{1}{2} \int \sqrt{u^{5/2}+u}
\d{u}.
$$
This is all correct, but leads to an integral far worse than the
original, especially if you notice that $x^2$  is ``almost'' the
derivative of $x^3+1.$  It is off by a factor of $3,$ but as we've seen,
constants can be brought outside of the integral, so this is not
really a problem.  So let's make $y=x^3+1,$  so that $\d{y}=3x^2 \d{x}$.  Thus 
$$
\int x^2 \sqrt{x^3+1} \d{x}=\int x^2 \sqrt{y} \left(\frac{\d{y}}{3x^2}\right)
  =\frac{1}{3} \int y^{1/2}  \d{y} =\frac{1}{3}
  \left(\frac{y^{\frac{3}{2}}}{\frac{3}{2}}\right)+ C=\frac{2}{9}
  (x^3+1)^{\frac{3}{2}}+C.
  $$

  So what's the definitive answer on a substitution?  Unfortunately,
  there is none.  Basically, you can try anything.  If it works,
  great!  If it doesn't, don't do something incorrectly to try and
  force it.  Try something else.  Something you can look for is parts
  of the integrand (the thing you are integrating) which are
  differentials of other parts of the integrand.  But there are no
  guarantees.  This is what makes integration harder than
  differentiation, much as long division is harder than
  multiplication.  But try something!  As you practice, you will gain
  more experience, so make sure you start looking for things that you
  have seen before.  And Practice, Practice, Practice!

  \begin{embeddedproblem}{}
    We've integrated the trigonometric functions $\sin(x), \cos(x),
    \tan(x), \cot(x).$  Now we will deal with secant and cosecant.
    \begin{enumerate}[label={\bf{}(\alph*)}]
    \item Consider the following ``trick:''
      $$
      \int \sec(x)  \d{x}=\int \sec(x)  \left(
        \frac{\sec(x)+\tan(x)}{\tan(x)+\sec(x)}\right)  \d{x}=\int
      \left(\frac{\sec^2(x)+\sec(x)  \tan(x)}{\tan(x)+\sec(x)}\right)  \d{x}.
        $$
        Do you recognize a substitution that can be used?  Use this to
        integrate $\sec(x).$ (Hint: Differentiate the denominator.)
      \item Use a trick similar to that in part a to compute $\int
        \csc(x)  \d{x}.$
    \end{enumerate}
  \end{embeddedproblem}

{\noindent{}\bf{}The Opposite of the Product Rule - Integration by
  Parts}
Substitution is a powerful tool, but it is not a ``cure all.''  For example, consider 
$$
\int x  \cos(x)  \d{x}.
$$

You can try, but there is no substitution that will really get you
anywhere.  Obviously, 
$$
\int x \cos(x)  \d{x}\ne\int x \d{x} \int \cos(x)  \d{x}
$$
so that is out of the question\footnote{Unfortunately this is a common
  error among beginners. We think it  is borne mostly of
  desparation. Integral Calculus is so filled with complex
  computations and arcane notation that it is easy to become
  overwhelmed and just want things to be easy. But that is just
  wishful thinking. Think about it. It is not true that if $h(x) =
  f(x)g(x)$ then $h^\prime(x) = f^\prime(x)g^\prime(x)$ -- we need the
  Product Rule to compute $f^\prime(x).$ What would make you think that the integral of a
  product would ever equal the product of the integrals?}.
Since the product rule is what we use to differentiate a product
elt's take a look at the Product Rule, with a view to how
we might run it backwards.   In general, we have the product rule as
$$
\d{uv}=u\d{v}+v\d{u}.
$$

Rearranging this a bit , we have
$$
u\d{v}=\d{(uv)}-v\d{u}
$$
$$
\int u \d{v}=\int (\d{(uv)}-v \d{u})=\int \d(uv) -\int v \d{u}=uv-\int
v \d{u}.
$$

This last formula is called the integration by parts formula. It can
be used to run the Product Rule backwards and is just what we need for
the previous integral.  Specifically, if we let $u=x$ and $\d{v}=\cos(x) \d{x},$
then $\d{u}=\d{x}$ and $v=\int \cos(x) \d{x}=\sin(x)$ (don't worry about the
arbitrary constant here for the moment).  Substituting these into our
integration by parts formula, $\int u \d{v}=uv-\int v \d{u},$ we have
$$
\int x  \cos(x)  \d{x}=x \sin(x)-\int \sin(x)  \d{x}=x \sin(x)-(-\cos(x)cb
)+C=x \sin(x) +  \cos(x)+C
$$

At this point, we should take note of a few things. First, notice that
this was not a substitution; we did not end up with and integral with
$u's$ or $v's$ in it.  These were introduced just to keep track of the
integration by parts formula.  Along that vein, notice that we needed
to use the entire integrand as opposed to a substitution where we
looked for parts of the integrand which were differentials of other
parts of the integrand.

Second, notice that like a substitution, the integration by parts
technique does not compute the integral.  Rather, it replaces that
integral with, hopefully, an easier integral.  This means that as long
as the integration by parts technique is applied correctly, it is not
wrong; it just might not be helpful.  For example, in  
$$
\int x \cos(x)  \d{x}
$$
we could have let $u=\cos(x)$ and $\d{v}=x \d{x}.$  This would give
$\d{u}=-\sin(x)  \d{x}$  and $v=\frac{x^2}{2}.$   Substituting these
into the integration by parts formula, we get 
$$
\int x  \cos(x)  \d{x}=\frac{x^22}  \cos(x)+\frac{1}{2} \int x^2
\sin(x)  \d{x}.
$$

Everything is correct, but we have an integral that is worse than
before.  So you will need to develop some experience with using this
technique as well.

Lastly, notice that we suppressed the arbitrary constant in $v.$  We
will see that this will not alter the result.

\begin{embeddedproblem}{}
\begin{enumerate}[label={\bf{}(\alph*)}]
\item 	In performing integration by parts on $\int x  \cos(x)  \d{x},$ suppose
  we let $u=x$ and $\d{v}=\cos(x)  \d{x}.$  As before, $\d{u}=\d{x},$ but suppose this time
  that we write $v=\sin(x)+a$ where $a$ is an arbitrary constant.  Show that
  we will obtain the same integral as before. 
\item 	Show that in general, if we let $\int \d{v}=v+a$ in the integration by
  parts formula, then we will obtain the same result as before, namely
   $ \int u \d{v}=uv-\int v \d{u}. $
 \end{enumerate}
\end{embeddedproblem}
\begin{embeddedproblem}{}
  	Go back and notice in our table of integrals that we had $\int e^x \d{x}=e^x+C,$ but we had no analog for $\ln(x).$  Perform integration by parts on 
$$
\int \ln(x)  \d{x}
$$
with $u=\ln(x)$ and $\d{v}=\d{x}$ to compute this integral.  (Note: Rather than trying to memorize this integral, just remember to use integration by parts.)
\end{embeddedproblem}

As with differentiation techniques and formulas, techniques such as
integration by parts and substitution are performed in concert,
utilizing whatever technique that is appropriate at the time.  As an
example of this consider

\begin{embeddedproblem}{}
  Perform integration by parts and then substitution to compute
$$
\int \inverse{sin}(x)  \d{x} \text{ and } \int \inverse{tan}(x) \d{x}.
$$

[We will apply integration by parts to $\int \inverse\sec(x)  \d{x}$ a
bit later, when we've developed some more techniques.] 
\end{embeddedproblem}

Just as integration by parts can be used with other techniques, it can
be used with itself multiple times.

\begin{embeddedproblem}{}
  Compute
  $$
  \int x^2e^x\d{x}.
  $$
  by using integration by parts twice.
\end{embeddedproblem}

As a word of advice, if you do perform integration by parts twice, be
sure not to switch the roles of $u$ and $v.$  This is not incorrect,
but will lead you back to where you started as indicated below 
$$
\int u \d{v}=uv-\int v \d{u}=uv-[vu-\int u \d{v}]=\int u \d{v}.
$$
Before we send you off to practice, we will look at one more ``trick''
one could do with integration by parts.  We will also illustrate a way
to keep track of $u, \d{v}, \d{u},$ and $ v$

\begin{myexample}
  Consider the integral $\int e^x  \sin(x)  \d{x}.$ If we let  $u=e^x$
  and  $\d{v}=\sin(x)  \d{x},$ then
  $$
  \begin{array}{cc}
    u=e^x & \d{v}=\sin(x)  \d{x}\\
    \d{u}=e^x  \d{x}&v=-\cos(x)    
  \end{array}
$$
so 
$$
\int e^x  \sin(x)  \d{x}=-e^x  \cos(x)+\int e^x  \cos(x)  \d{x}
$$
\end{myexample}

We can apply integration by parts again, being careful not to ``switch
the roles of $u$ and $v.$''
$$
\begin{array}{cc}
  u=e^x& \d{v}=\cos(x)  \d{x}\\
\d{u}=e^x  \d{x} & v=\sin(x)
\end{array}
$$
so 
$$
\int e^x  \sin(x)  \d{x}=-e^x  \cos(x)+\int e^x  \cos(x)  \d{x}
$$
$$
\int e^x  \sin(x)  \d{x}=-e^x  \cos(x)+[e^x  \sin(x)-\int e^x  sin(x)
\d{x}]=-e^x  \cos(x)+e^x  \sin(x)-\int e^x  \sin(x)  \d{x}.
$$

It looks like we went full circle and came back to where we started.  However, if we let $I=\int e^x  \sin(x)  \d{x},$ we really have the equation
$$
I=-e^x  \cos(x)+e^x  \sin(x) -I
$$
and we can solve for $I$ (remembering to put in our arbitrary constant
$c$).
$$
2I=-e^x  \cos(x)+e^x  \sin(x)
$$
$$
\int e^x  (x)sin(x)  \d{x}=I=-e^x  \cos(x)+e^x  \frac{\sin(x)}{2}+C
$$
So basically, when you are integrating you can use any mathematically
legal trick you wish.  As long as you do that, what you are doing is
not wrong.  It may not help, but it is not wrong.  Again experience
helps you to recognize fruitful paths to take.  So again, Practice,
Practice, Practice!

\begin{PracticeProblemSection}
  
\end{PracticeProblemSection}

{\noindent{}\bf{}Trigonometric Substitutions -- Why Simplify When We
  Can Complicate?!}
Recall that earlier in this section we encountered the integral:
$$
\int \frac{1}{\sqrt{1+z^2}}  \d{z}
$$
in trying to determine the solution to the catenary.  We actually have
the tools to solve this type of integral once we notice that
(surprisingly!) we would like to use the substitution $z=\tan(\theta)$.  As
strange as it sounds to take a perfectly good integral and purposely
insert trigonometry into it, this is just the ticket here.  Since we
have the trigonometric identity $\tan^2(\theta)+1=\sec^2(\theta),$ then this
substitution is just the thing to rid us of the square root, which
poses more of an issue than the trigonometry does.  Specifically, if
we let $z=\tan(\theta),$ then $\d{z}=\sec^2(\theta)$  $\d{\theta},$ and so we have 
$$
\int \frac{1}{\sqrt{1+z^2}}  \d{z}=\int \frac{\sec^2(\theta)}{\sqrt{1+\tan^2(\theta)}}  \d{\theta}=\int \frac{\sec^2(\theta)}{\sec(\theta)}   \d{\theta}=\int \sec(\theta) \d{\theta}=\ln(\tan(\theta)+\sec(\theta)+C
$$
The trick now is to put things back in terms of $z.$  We have
$\tan(\theta)=z,$ but what about $\sec(\theta{})?$  We could certainly
utilize the trigonometric 
identity from before, but a more general way to do this is to draw a
triangle.  The beauty of this is that we don't know what the angle $\theta$ 
is, nor do we really care.  All we need is that $\tan(\theta)=\frac{z}{1}.$  This leads
to the following triangle
\includegraphics*[height=2in,width=2in]{../Figures/TrigSub1}

With this triangle, not only can we obtain the aforementioned $\tan(\theta)=\frac{z}{1}$ and $\sec(\theta)=\frac{\sqrt{1+z^2}}{1}$ which is what we need in the problem, but we could also obtain the following if we needed it.
$$
\sin(\theta)=\frac{z}{\sqrt{1+z^2}}, \cos(\theta)=\frac{1}{\sqrt{1+z^2}},
\cot(\theta)=\frac{1}{z}, \csc(\theta)=\frac{\sqrt{1+z^2}}{z}
$$
We really urge you to do this (DRAW A TRIANGLE) instead of relying on
memorizing trigonometric identities which actually are derived from
the triangle.  This gives us that the differential equation 
$$
\int \frac{1}{\sqrt{1+z^2}},  \d{z}=\int \frac{w}{H}  \d{x}
$$
is satisfied by
$$
\ln(z+\sqrt{1+z^2})=\frac{w}{H} x+C.
$$

\begin{embeddedproblem}{}
  \begin{enumerate}[label={\bf{}(\alph*)}]
  \item Utilize that $z=\dfdx{y}{x}$ for the catenary and that  $z(0)= \left.\dfdx{y}{x}\right|_{x=0}=0$ (this was the low point on the hanging chain) to show that $c=0$ and so
$$
\dfdx{y}{x}=z=\frac{1}{2} \left(e^{\frac{w}{H} x}-e^{-\frac{w}{H} x}\right).
$$
\item 	Integrate the result in part (a) to obtain 
$$
y=\frac{H}{w} \left(\frac{e^{\frac{wx}{H}}+e^{-\frac{wx}{H}}}{2})\right)+c
$$
which is the equation of the catenary as stated earlier.
\end{enumerate}
\end{embeddedproblem}

Recall Problem \#52 from Section 3.2.

Consider the top view of a plane $P$ starting at the point $(1,0)$ and
traveling up the line $x=1$ at a constant speed $v.$  When the plane is at
$(1,0),$ a homing missile $M$ is fired from the origin directly at the
plane. Assuming that the missile travels at a speed which is $k$ times
the speed of the plane $(k>1)$ and is always aimed directly at the
plane, we will find the path the missile takes.  The diagram below
shows the situation at time $t.$
\includegraphics*[height=2in,width=2in]{../Figures/PursuitCurve}

If we let s denote the distance the missile has traveled at time $t,$
show that the missile's path must satisfy the IVP:
$$
\dfdx{y}{x}=\frac{\frac{s}{k}-y}{1-x},   y(0).
$$

Problem \#52 had you then apply Euler's Method to approximate this
pursuit curve that the missile follows.  We didn't give you an
equation for the curve, but now we actually have the means to find
it.
\begin{embeddedproblem}{}
  \begin{enumerate}[label={\bf{}(\alph*)}]
  \item Differentiate the above equation to show that the missile's
    path must satisfy the differential equation:
    $$
    (1-x)\dfdxn{y}{x}{2}=\frac{1}{k}\dfdx{s}{x}=\frac{1}{k} \sqrt{1+\left(\dfdx{y}{x}\right)^2 }
    $$
    with the initial conditions $y(0)=0,$ and $\dfdx{y}{x}(0)=0.$
  \item	Let $z=\dfdx{y}{x},$  so that the differential equation in
    part (a) becomes
    $$
    \frac{1}{\sqrt{1+z^2}}  \d{z}=\frac{1}{k}\cdot\frac{1}{1-x}, \ \
    z(0)=0
    $$
    Utilize the trigonometric substitution we used in the solution of
    the catenary to show that:
    $$
    z=\frac{1}{2} \left[(1-x)^{-\frac{1}{k}}-(1-x)^{\frac{1}{k}}
    \right].
    $$
  \item Use the fact that $z=\dfdx{y}{x}$ and $y(0)=0$ to show that
    $$
    y=\frac{k}{2}\left[\frac{(1-x)^{\frac{k+1}{k}}}{k+1}-\frac{(1-x)^{\frac{k-1}{k}}}{k-1}\right]+\frac{k}{k^2-1}
%    y=\frac{k}{2}\left[\frac{(1-x)^{\frac{k+1}{k}}}{k+1}-\frac{(1-x)^{\frac{k-1}{k}}{k-1}}\right]+\frac{k}{k^2-1}.
    $$
  \item How far has the plane gone when the missile reaches it?  What
    happens as $k\rightarrow{}1?$
  \end{enumerate}
\end{embeddedproblem} 

Of course, the tangent function is not the only trigonometric function
at our disposal.  Recall the Tractrix Problem (Problem in Context
\#$51$ in Section $3.2$).

The following is the top view of a tractor-trailer.  Initially, the
center of the rear axle of the tractor is at the origin and the center
of the rear axle of the trailer is at the point $(1,0).$
\centerline{
  \includegraphics*[height=3in,width=5in]{../Figures/Tractrix}}
Suppose the tractor pulls the front wheels up the y-axis and that the
rear wheels don't slip.  In the problem you had to show the path
$y=y(x)$ that the center of the rear axle of the trailer follows must
satisfy the equations 
$$
\dfdx{y}{x}=-\frac{\sqrt{1-x^2}}{x}, \ \ y(1)=0.
$$

You then utilized Euler's Method to approximate the curve that the
rear wheels follow (the tractrix).  We are now in a position to solve
this equation.
\begin{embeddedproblem}{}
  \begin{enumerate}[label={\bf{}(\alph*)}]
  \item Find the equation $y=y(x)$ for the tractrix.  [Hint:  After you
    separate the variables, you can substitute $x=sin(\theta)$ in your
    integral.]   
  \item Graph your answer and see how it compares to the graph plotted
    back in Problem $51$ in Section\footnote{\textcolor{red}{Bud - Do
        we want to include the plot obtained in Problem 51?}} $3.2.$
  \end{enumerate}
\end{embeddedproblem}

Again, having trigonometry involved was a better option than having
the square root.  Here, we utilized the trigonometric identity
$\sin^2(\theta)+\cos^2(\theta)=1.$  A surprising number of integrals involve terms such
as $\sqrt{1-x^2},$ $\sqrt{1+x^2},$ or $\sqrt{x^2-1}.$  In these cases, it is often
advantageous to utilize a trigonometric substitution, using the
identities $\cos(\theta)=\sqrt{1-\sin^2(\theta)},$ $\sec(\theta)=\sqrt{1+\tan^2(\theta)},$ or $\tan(\theta)=\sqrt{\sec^2(\theta)-1}$ to
remove the square root, which typically is more a concern than the
trigonometry.  However, you should remember that there are other
techniques at your disposal, so you need to be judicious about what
you utilize.

\begin{embeddedproblem}{}
  For example, consider $\displaystyle\int \frac{x}{\sqrt{1-x^2}}\d{x}.$
  \begin{enumerate}[label={\bf{}(\alph*)}]
  \item Utilize a trigonometric substitution to compute the above
    integral.
  \item Utilize a non-trigonometric substitution to compute the above
    integral.
  \item Which did you find more appealing to do?
  \end{enumerate}
\end{embeddedproblem}
\begin{embeddedproblem}{}
As you recall from our original integration table, we had \\
%  \begin{wraptable}{r}{5in}
  \begin{tabular}{ |c|c| }
    \hline
    % \multicolumn{2}{|c|}{\bf{}General} \\
    % \multicolumn{2}{|c|}{\bf{}Integration Rules} \\\hline
    % Differentiation Rule & Integration Rule\\\hline
    $\d(\inverse{\sin}(x) )=\frac{1}{\sqrt{1-x^2}}  \d{x}$& $\int\frac{1}{\sqrt{1-x^2}}\d{x}=\inverse{\sin}(x)+c$\\\hline
    $\d(\inverse{\tan}(x))=\frac{1}{1+x^2}  \d{x}$&$\int\frac{1}{1+x^2}  \d{x}=\inverse{\tan}(x)+C$\\\hline
    $\d(\inverse{\sec}(x))=\frac{1}{x\sqrt{x^2-1}}  \d{x}$&$\int \frac{1}{x\sqrt{x^2-1}}  \d{x}=\inverse{\sec}(x)+C$\\\hline{}
  \end{tabular}
%\end{wraptable}
In case you forget these integrals, don't worry, you can use a
trigonometric substitution to obtain them.
\begin{enumerate}[label={\bf{}(\alph*)}]
 \item 	Use the substitution $x=sin(\theta)$ to compute 
   $\int \frac{1}{\sqrt{1-x^2}}  \d{x}.$
 \item Use the substitution $x=\tan(\theta)$ to compute 
 $\int\frac{1}{\sqrt{1+x^2}} \d{x}$
 \item 	Use the substitution $x=\sec(\theta)$ to compute 
   $\int \frac{1}{x\sqrt{x^2-1}}  \d{x}.$
\end{enumerate}
\end{embeddedproblem}

As we mentioned, trigonometry can be used for integrals involving $\sqrt{1-x^2},\sqrt{1+x^2},$ or $\sqrt{x^2-1}.$  What if we had
$$
\int \frac{\d{x}}{\sqrt{4+x^2}}?
$$


Rather than learning a set of rules, we can utilize some algebra to
transform it into something we know.  The trick is that we really want
a $1$ where the $4$ is.  We can't just change it, but we can factor it
out.
$$
\int \frac{\d{x}}{\sqrt{4+x^2}} = \int
\frac{\d{x}}{\sqrt{4\left(1+\frac{x^2}{4}\right)}}=\frac{1}{2} \int
\frac{\d{x}}{\sqrt{\left(1+\left(\frac{x}{2}\right)^2\right)}}.
$$

Now we can let  $\frac{x}{2}=\tan(x),$ and proceed as before.  
Similarly,
$$
\int \frac{\d{x}}{\sqrt{5-x^2}}=\int
\frac{\d{x}}{\sqrt{5(1-x^2/5)}}=\frac{1}{\sqrt{5}} \int
\frac{\d{x}}{\sqrt{1-\left(\frac{x}{\sqrt{5}}\right)^2}},
$$
and we can let $\frac{x}{\sqrt{5}}=\sin(x).$

We can even handle something like
$$
\int \frac{\d{x}}{\sqrt{x^2-6x}}
$$
by completing the square first.
$$
\int \frac{\d{x}}{\sqrt{x^2-6x}}=\int
\frac{\d{x}}{\sqrt{x^2-6x+9-9}}=\int
\frac{\d{x}}{\sqrt{(x-3)^2-9}}=\frac{1}{3} \int
\frac{dx}{\sqrt{\left(\frac{x-3}{3}\right)^2-1}}
$$

Now we can let $\frac{x-3}{3}=\sec(x).$  As long as the expression
inside the radical is a quadratic, we can manipulate it algebraically
to produce something of the form
$\sqrt{1+\left(\frac{x-b}{a}\right)^2},$
$\sqrt{1-\left(\frac{x-b}{a}\right)^2},$  or
$\sqrt{\left(\frac{x-b}{a}\right)^2-1}.$
From there, we can use the appropriate
trigonometric substitution.

\begin{PracticeProblemSection}
  
\end{PracticeProblemSection}

The trigonometric substitutions we've dealt with so far have dealt
with square roots.  How about something like
$$
\int\frac{\d{x}}{\left(1-x^2\right)^{\frac32}?}
$$
If we utilize the substitution $x=\sin(\theta{}),$ then $\d{x}=\cos(\theta) \d{\theta},$ so
$$
\int \frac{\d{x}}{\left(1-x^2\right)^{\frac{3}{2}}} =\int
  \frac{\cos(\theta)}{\left(1-\sin^2(\theta)\right)^3} \d{\theta}=\int \frac{\cos(\theta)}{\cos^3(\theta)}
    \d{\theta}=\int \sec^2(\theta)  \d{\theta}=\tan(\theta)+C.
    $$

Utilizing the fact that $\sin(\theta)=\frac{x}{1},$ we get the
triangle:\\
\centerline{\includegraphics*[height=2in,width=2in]{../Figures/TrigSub2}}
so
$$
\int \frac{\d{x}}{(1-x^2 )^{\frac{3}{2}}}
=\tan(\theta)+C=\frac{x}{\sqrt{1-x^2}}+C.
$$

Fortunately, we were still able to handle this, but what if we started with the integral
$$
\int\frac{\d{x}}{(x^2-1)^{\frac{3}{2}}}?
$$
Utilizing the substitution $x=\sec(\theta)$ would lead to the integral
$$
\int\frac{\sec(\theta)}{\tan^2(\theta)}   \d{\theta}.
$$

What do we do with this one?  We can rewrite this as
$$
\int\frac{\sec(\theta)}{\tan^2(\theta)}   \d{\theta}=\int
\frac{\frac{1}{\cos(\theta)}}{\frac{\sin^2(\theta)}{\cos^2(\theta)}}
\d{\theta} =\int\frac{\cos(\theta)}{\sin^2(\theta)}  \d{\theta}=
\int\cot(\theta)  \csc(\theta)  \d{\theta}=-\csc(\theta)+C
$$

This underscores the need to integrate a number of integrals involving
trigonometric functions.  Again, it is a trial and error kind of
thing, but there are some standard tricks to learn.  It is probably
best to learn these tricks by looking at some specific examples, but
as you do, focus on the technique being used and how it can be
utilized on similar integrals.

\begin{myexample}
  $\displaystyle \int\sin^3(\theta)\cos^2(\theta)\d{\theta}.$\\
  Since sine is to an odd power, it makes sense to ``save one of the
  sines'' and convert everything else to cosine, utilizing the identity
  $\sin^2(\theta)=1-\cos^2(\theta),$ and then applying a substitution. 
  $$
  \int\sin^3(\theta)  \cos^2(\theta)  \d{\theta}=\int\sin(\theta)
  \sin^2(\theta)  \cos^2(\theta)
  \d{\theta}=\int(1-\cos^2(\theta))\cos^2(\theta)  \sin(\theta)
  \d{\theta}.
  $$
  
Letting $z=\cos(\theta),$ $\d{z}=-\sin(\theta)  \d{\theta},$ so we have
$$
\int\sin^3(\theta)  \cos^2(\theta)  \d{\theta}=\int(1-z^2) z^2
(-\d{z})=\int(z^4-z^2
)\d{z}=\frac{z^5}{5}-\frac{z^3}{3}+C=\frac{\cos^5(\theta)}{5}-\frac{\cos^3(\theta)}{3}+C
$$
\end{myexample}
Notice that if cosine had been to an odd power, you could have done
the same trick, focusing on saving one cosine as part of the
differential, and converting the rest to sine.  Notice also that there
would have been nothing ``wrong'' with doing the following.
$$
\int\sin^3(\theta)  \cos^2(\theta)  \d{\theta}=\int\sin^3(\theta)
\cos(\theta)  \cos(\theta)
\d{\theta}=\int\sin^3(\theta)\sqrt{1-\sin^2(\theta)}   \cos(\theta)
\d{\theta}.
$$

Letting $x=\sin(\theta),$ $\d{x}=\cos(\theta)\d{\theta}$ would give
$$
\int\sin^3(\theta) \cos^2(\theta) \d{\theta}=\int x^3\sqrt{1-x^2 }
\d{x}
$$

which is potentially a harder integral than the original (though you
could do it by parts).  Remember that we introduced trigonometry to
handle such integrals.  

\begin{myexample}
  
  What happens if you don't have an odd power of sine or cosine?  For
  example, suppose you have 
$$
\int\sin^2(\theta)  \cos^2(\theta)  \d{\theta}.
$$

Saving a sine or cosine would lead to an integrand with a square root,
which we are trying to avoid, if we can help it.  In this case, there
are a couple of trigonometric identities that can help 
$$
\sin^2(\theta)=\frac{1}{2} [1-\cos(2\theta)]
\cos^2(\theta)=\frac{1}{2} [1+\cos(2\theta) ]
$$

Utilizing these, we have
$$
\int\sin^3(\theta)  \cos^2(\theta)  \d{\theta}=\frac{1}{4}
\int(1-\cos(2\theta) )(1+\cos(2\theta) )\d{\theta}=\frac{1}{4}
\int(1-\cos^2(2\theta) )  \d{\theta}=\frac{1}{4} \int\sin^2(2\theta)
\d{\theta}
$$

Utilizing this ``half angle'' identity again, we obtain
$$
\frac{1}{4} \int\sin^2(2\theta)  \d{\theta}=\frac{1}{8}
\int(1-\cos(4\theta) )\d{\theta}=\frac{1}{8} \left(\theta-\frac{1}{4}
  \sin(4\theta) \right)+C
$$
\end{myexample}
\begin{embeddedproblem}{}
  Show that the same result would have occurred if we used the double angle identity
$$
\sin(2\theta)=2 \sin(\theta)  \cos(\theta).
$$
\end{embeddedproblem}
\begin{myexample}
  $$
  \int \sec(4\theta)  \tan(2\theta)  \d{\theta}
  $$
  In this example, we can save a $\sec(2\theta)$ as part of the differential
  and utilize the identity $\sec^2(\theta)=1+\tan^2(\theta)$ to change everything else
  into an expression in tangent.  Since the secant is to an even power
  this again avoids square roots, which is the general idea. 
$$
\int \sec(4\theta)  \tan^2(\theta)  \d{\theta}=\int \sec^2(\theta)
\tan^2(\theta)  \sec^2(\theta) \d{\theta}=\int
(1+\tan^2(\theta))\tan^2(\theta)  \sec^2(\theta) \d{\theta}
$$

Letting $y=\tan(\theta),$ we have $\d{y}=\sec^2(\theta)  \d{\theta},$ so we obtain
$$
\int (1+y^2)y^2\d{y}=\int (y^2+y^4
)\d{y}=\frac{y^3}{3}+\frac{y^5}{5}+C=\frac{\tan^3(\theta)}{3}+\frac{\tan^5(\theta)}{5}+C 
$$
\end{myexample}

\begin{myexample}
  If tangent occurs to an odd power, we can save a tangent and a
  secant as part of the differential and use the same identity to
  convert everything else to secant. 
$$
\int \sec^3(\theta)  \tan^3(\theta)  \d{\theta}=\int \sec^2(\theta)
\tan^2(\theta)  \sec(\theta)  \tan(\theta) \d{\theta}=\int
\sec^2(\theta) (\sec^2(\theta)-1)  \sec(\theta)  \tan(\theta)
\d{\theta}
$$

Letting $u=\sec(\theta),$ $ \d{u}=\sec(\theta)  \tan(\theta)
\d{x}{\theta},$ so we have 
$$
\int u^2 (u^2-1)\d{u}=\int (u^4-u^2
)\d{u}=\frac{u^5}{5}-\frac{u^3}{3}+C=\frac{\sec^5(\theta)}{5}-\frac{\sec^3(\theta)}{3}+C
$$
\end{myexample}
\begin{embeddedproblem}{}
  Consider 
  $\displaystyle\int\sec^4(\theta)  \tan^3(\theta)  \d{\theta}.$
\begin{enumerate}[label={\bf{}(\alph*)}]
\item 	Compute this integral by saving a $\sec^2(\theta)$ and converting the
  rest into an expression in tangent. 
\item 	Compute this integral by saving a $\sec(\theta)  \tan(\theta)$ and converting
  the rest into an expression in secant. 
\item 	Verify that the two answers you obtained are really the same. 
\end{enumerate}
\end{embeddedproblem}

\begin{embeddedproblem}{}
  Previously, we computed $\int \tan(\theta) \d{\theta}$ by rewriting
  $\frac{\tan(\theta)\sin(\theta)}{\cos(\theta)}.$ Alternatively,
  compute this integral by rewriting it as
  $$
  \int \tan(\theta) \d{\theta}=\int \frac{\sec(\theta)
    \tan(\theta))}{\sec(\theta)} \d{\theta}.
  $$

  The same tricks we applied with secant and tangent would apply
  equally well to integrals involving cosecant and cotangent utilizing
  the identity $\csc^2(\theta)=1+\cot^2(\theta).$ Of course, we have
  not covered every type of integral that can occur and you can always
  find integrals that will take other tricks.  We've only provided
  some basics that you can build on as you gain experience.  So
  practice, practice, practice!
\end{embeddedproblem}

\begin{PracticeProblemSection}
  \begin{enumerate}
  \item Trig Integrals
  \item More trig substitutions
  \end{enumerate}
\end{PracticeProblemSection}

{\noindent{}\bf{}Back to Logistic Growth: Partia Fractions}

Recall that in the graphing section of this book (Section 6.4) we had
an example for logistic growth 
$$
\dfdx{y}{t}=.01y(1-\frac{y}{100})=.0001y(100-y)
$$

In that chapter, we analyzed this growth rate to draw a qualitative
graph of $y=y(t),$ but we really could not produce specific values of
$t$ for where things happened, because we dind't have a formula for
$y(t).$  Separating the variables is no problem
$$
\dfdx{y}{100-y} =.0001 \d{t}.
$$

To integrate the left-hand side, we could expand the denominator to
complete the square and apply the appropriate trigonometric
substitution.

As we've mentioned before, all mathematical models are simplifications
of reality, and typically are modified to take into account more
complexities.  For example, consider a modification of the logistic
model above 
$$
\dfdx{y}{t}=.01y\left(1-\frac{y}{100}\right)\left(\frac{y}{10}-1\right)=.00001y(100-y)(y-10)
$$
In this case, 10 is called the \emph{minimum viability level of the
population.} (Why?)  Solving this would require solving
$$
\int \frac{\d{y}}{(y(100-y)(y-10))}
$$
which is not really conducive to a trigonometric substitution.

Here, we will examine a technique which does not involve trigonometry
and is really algebraic in nature.  It will also be useful for
examining integrals such as the two above.  The trick to try and
separate the fraction into ``partial fractions'' which will be easier to
integrate.  For example, consider $\frac{1}{y(100-y)}.$  Unfortunately, we
cannot just expand the denominator and divide this into two fractions
as we could if it was the numerator.  This is due to the fact that
fractions are added by finding a common denominator, not just adding
denominators.  However it is reasonable to make an educated guess that
such a fraction can be written as
$$
\frac{1}{y(100-y)} =\frac{A}{y}+\frac{B}{100-y}
$$
for some as yet to be determined $A$ and $B.$  We can see if this
guess pans out by combining the partial fractions to obtain the
original. 
$$
\frac{1}{y(100-y)}
=\frac{A}{y}+\frac{B}{100-y}=\frac{A(100-y)+By}{y(100-y)}
=\frac{100A+(B-A)y}{y}.
$$
This says that 
\begin{align*}
B-A&=0
100A&=1
\end{align*}
so that $A=\frac{1}{100}$ and $B=A=\frac{1}{100}.$  Hence 
$$
\frac{1}{y(100-y)} =\frac{1/100}{y}+\frac{1/100}{100-y}.
$$
Back to our integration problem, we would obtain
$$
\int \frac{\d{y}}{y(100-y)} =\int \frac{1/100}{y}\frac{+1/100}{100-y}
\d{y}=\frac{1}{100} [\int \frac{\d{y}}{y}+ \int
\frac{\d{y}}{100-y}=\frac{1}{100} [\ln{y}-\ln(100-y)]+C.
$$

\begin{embeddedproblem}{}
\begin{enumerate}[label={\bf{}(\alph*)}]
\item 	Use the integral provided above to obtain an equation for $y$
  in the logistic growth model 
  $$
  \dfdx{y}{t}=.01y\left(1-\frac{y}{100}\right),   y(0)=y_0
  $$
  for the case where $0<y<100.$ 
\item 	Solve the same problem for the case where y>100.  [This is a
  spot where you need to look at the absolute value in the logarithm.
  We told you it comes up occasionally.]
\end{enumerate}
\end{embeddedproblem}

Now that we've shown you an example where decomposing a fraction of
polynomials (called a rational function) into a sum of partial
fractions is practical, let's look into a systematic way of doing this
algebraically.  All of this hinges on the following theoretical fact
about polynomials.

Suppose we have two polynomials p(x) and q(x) which have no common
factors (other than constants, which are considered trivial.  For
example, $p(x)=2\left(\frac{p(x)}{2}\right),q(x)=2\left(\frac{q(x)}{2}\right),$ which would make $2$ a trivial
common factor of $p(x)$ and $q(x).$)  Then there are polynomials $r(x)$ and
$s(x)$ such that
$$
r(x)\cdot p(x)+s(x)\cdot q(x)=1.
$$
In other words, there is some (linear) combination of $p$ and $q$ which produces $1.$  For example, notice that
$$ 
(x+a)-(x+b)=a-b
$$	
so if $a\ne b,$ then
$$
\frac{1}{a-b} (x+a)+\frac{-1}{a-b} (x+b)=1
$$

We determined this combination in a somewhat ad-hoc way, but there is
a systematic way to do this involving no more than long division of
polynomials.  We won't delve into this, but will adopt a more
(educated) guess and check method.  Notice that in the above example 
\begin{align*}
\frac{1}{(x+a)(x+b)} &=\frac{\frac{1}{a-b} (x+a)+\frac{-1}{a-b}
  (x+b)}{(x+a)(x+b)}\\
                      &=\frac{\frac{1}{a-b}}{x+b}+\frac{\frac{-1}{a-b}}{x+a}\\
                      &=\frac{\frac{-1}{a-b}}{x+a}+\frac{\frac{1}{a-b}}{x+b}.
\end{align*}

In general, if we have
$$
r(x)\cdot p(x)+s(x)\cdot q(x)=1
$$
then
$$
\frac{1}{p(x)\cdot q(x)}=\frac{r(x)\cdot p(x)+s(x)\cdot
  q(x)}{p(x)\cdot q(x)}=\frac{s(x)}{p(x)} +\frac{r(x)}{q(x)}.
$$

Notice that if $t(x)$ is any polynomial, then
$$
\frac{t(x)}{p(x)\cdot q(x)}=\frac{t(x)\cdot s(x)}{p(x)}
+\frac{t(x)\cdot r(x)}{q(x)}
=\frac{A(x)}{p(x)} +\frac{B(x)}{q(x)}
$$
where $A(x)=t(x)\cdot s(x)$ and $B(x)=t(x)\cdot r(x)$ are two polynomials.

To summarize all of this, we have, in theory, that if $p(x)$ and
$q(x)$ are two polynomials with no non-trivial common factors, then
for any polynomial $t(x),$ there are polynomials $A(x)$ and $B(x)$
such that 
$$
\frac{t(x)}{p(x)\cdot q(x)}=\frac{A(x)}{p(x)} +\frac{B(x)}{q(x)}.
$$
To extend this idea further, if $p(x), q(x),$ and $r(x)$ have no
nontrivial factors, then there are polynomials $A(x),B(x),$ and $C(x)$ with
$$
\frac{t(x)}{p(x)\cdot{}q(x)\cdot{}r(x)}=\frac{A(x)}{p(x)}
+\frac{B(x)}{q(x)} +\frac{C(x)}{r(x)}
$$
and this can be extended to any number of factors in the denominator.
Again, this can be proven theoretically, and this partial fractions
decomposition can be obtained systematically, but we will adopt a
guess and check method to find $A, B,$ and $C.$ To make our guess a
little more educated, the following fact can be employed.  You can
make peace with this in your own closet.

If the degree of $t(x)$ is less than the degree of $p(x)\cdot
q(x)\cdot r(x),$ then $A,B,$ and $C$ can be chosen with
$\deg(A)<\deg(p),$ $\deg(B)<\deg(q),$ and $\deg(C)<\deg(r).$

\begin{myexample}
  Recall we had $\displaystyle\int \frac{\d{y}}{y(100-y)(y-10)}.$ To
  compute this we will find the partial fractions decomposition of
  $\frac{1}{y(100-y)(y-10)}.$ We make the educated guess
  $$
  \frac{1}{y(100-y)(y-10)} =\frac{A}{y}+\frac{B}{100-y}+\frac{C}{y-10}
  $$
  and determine what $A,B,$ and $C$ are.  This can be done in a number
  of ways, but the most direct (and labor intensive) way to do this is
  to combine the terms in the right-hand side and compare
  coefficients.
  \begin{align*}
    \frac{1}{y(100-y)(y-10)}
    &=\frac{A(100-y)(y-10)+By(y-10)+Cy(100-y)}{y(100-y)(y-10)}\\
    &=\frac{A(-y^2+110y-100)+B(y^2-10y)+C(100y-y^2
      )}{y(100-y)(y-10)}\\
    &=\frac{(-A+B-C) y^2+(110A-10B+100C)y+(-1000A)}{y(100-y)(y-10)}.
  \end{align*}
Since the fractions must be equal and the denominators are the same,
then we can set the numerators equal. 
$
1=0y^2+0y+1=(-A+B-C) y^2+(110A-10B+100C)y+(-1000A)
$
\begin{align*}
-A+B-C&=0
110A-10B+100C&=0
-1000A&=1
\end{align*}
Solve this anyway you wish, but you should get
$$
A=-1/1000,   B=1/9000,   C=1/900.
$$
Thus
$$
\frac{1}{y(100-y)(y-10)} =\frac{-1/1000}{y}+\frac{1/9000}{100-y}+\frac{1/900}{/y-10}
$$
so
\begin{align*}
\int \frac{\d{y}}{y(100-y)(y-10)} &=\frac{-1}{1000} \int
\frac{\d{y}}{y+1/9000} \int \frac{\d{y}}{100-y}+\frac{1}{900} \int \frac{\d{y}}{y-10}
&=-1/1000  \ln(y)-1/9000  \ln(100-y)+1/900  \ln(y-10)+C.
\end{align*}
\end{myexample}

\noindent{\bf{}A few Practice Problems with linear terms and no shortcuts.}\\
The practice problems above illustrate that determining the partial
fractions decomposition is really algebra and not calculus.  Bear in
mind that the algebra is tedious, but it is algebra.  To eliminate
some of the tedium, there are some shortcuts.  We will explore one of
them here.  We had you practice some without the shortcut because one
should not follow a shortcut unless he/she is familiar with the
non-shortcut.  Furthermore, shortcuts don't always pan out; the
traditional algebraic method will always work.  We'll illustrate this
with an example.

\begin{myexample}
  Suppose we wish to find the partial fractions decomposition of
$\frac{3x+1}{((x-1)(x-2)(x-5))}.$  As before, we have
\begin{equation}
  \label{eq:partfrac1}
  \frac{3x+1}{(x-1)(x-2)(x-5)}=
  \frac{A}{x-1}+\frac{B}{x-2}+\frac{C}{x-5}
\end{equation}
               
This leads to
$$
3x+1=\frac{A}{(x-2)(x-5)}+\frac{B}{(x-1)(x-5)}+\frac{C}{(x-1)(x-2)}
$$
Observing that
\begin{enumerate}
\item this formula will be true for every real (or complex) value of
  $x$ and that,
\item the factor $x-1$ appears in two of the terms on the right
\end{enumerate}
suggests that we let $x=1.$ In that case we have
$$
3(1)+1=A(1-2)(1-5)
$$
or
$$
(3(1)+1)/(1-2)(1-5) =A.
$$

We did not simplify this last formula in order to highlight the
following. Observe that we get the same result if we simply ``cover
up'' the factor $(x-1)$ on the left side of
equation~\ref{eq:partfrac1} above and set $x=1.$ Then
$$
A=\frac{3(1)+1}{(covered)(1-2)(1-5)}=1.
$$

In precisely the same fashion (setting $x=2$ and $x=3$) we get:
$$
B=\frac{(3)2+1}{(2-1)(covered)(2-5)} =  \frac{-7}{3},
$$
and
$$
C=\frac{(3)5+1}{(5-1)(5-2)(covered)}=\frac{4}{3},
$$
so that
$$
\frac{3x+1}{(x-1)(x-2)(x-5)}=\frac{1}{x-1}+\frac{\frac{-7}{3}}{x-2}+\frac{4}{3}{x-5}. 
$$
This trick is known as "Heaviside's Cover-Up Method" (HCUM). It is
well known, and can be used to find the PFD of a rational function as
long as the denominator can be factored into distinct linear factors.

\begin{PracticeProblemSection}
  Insert problems on HCU
\end{PracticeProblemSection}

\noindent{\bf{}Expanding Partial Fractions Decompositions}\\
In the previous partial fractions decompositions, we started with a
rational function where the degree of the numerator was less than the
degree of the denominator and where the factors in the denominator
were linear.  In this section, we will expand our techniques to
encompass other situations.

For example, suppose the degree of the numerator is not less than the
degree of the denominator.  If this is the case, then we can perform
long division first and then find the partial fractions decomposition
of the remainder.  For example, suppose we have 
$$
\frac{x^3-2x+1}{x^2+x-6}=\frac{x^3-2x+1}{(x-2)(x+3)}.
$$
We first perform long division of the original rational function.
$$
\frac{x^3-2x+1}{x^2+x-6}=x-1+\frac{5x-5}{x^2+x-6}=x-1+5\left(\frac{x-1}{(x-2)(x+3)}\right).
$$
We can now perform a partial fractions decomposition on the remaining fraction.

\begin{PracticeProblemSection}
  INsert practice exercises with integrals with involving rational
  functions requiring long division. 
\end{PracticeProblemSection}

What about something like
$$
\frac{3x+2}{x^2 (x-1) }
$$

Since $x^2$ and $x-1$ have no common (nontrivial) factors, then we can
still preform our partial fractions decomposition.  Recall that in the
decomposition, the numerators in the individual terms can be chosen so
that the degrees are less than the degrees in the denominators.  This
being said, we make the following educated guess as to the form of the
partial fractions decomposition.
$$
\frac{3x+2}{x^2 (x-1)}=\frac{Ax+B}{x^2} +\frac{C}{x-1}.
$$
We can now recombine this to determine $A, B,$ and $C.$ 
\begin{align*}
\frac{3x+2}{x^2 (x-1)}&=\frac{Ax+B}{x^2} +\frac{C}{x-1}\\
                      &=\frac{(Ax+B)(x-1)+Cx^2}{x^2 (x-1)}\\
                      &=\frac{Ax^2+Bx-Ax-B+Cx^2}{x^2 (x-1)}\\
                      &=\frac{(A+C) x^2+(-A+B)x-B}{x^2 (x-1)}
\end{align*}
This leads to
\begin{align*}
A+C&=0\\
-A+B&=3\\
-B&=2
\end{align*}
Solving this (any way you can) we get $A=-5, B=-2,$ and $C=5$ so that
$$
\frac{3x+2}{x^2 (x-1)}=\frac{-5x-2}{x^2} +\frac{5}{x-1}.
$$
If we wanted to integrate this, we would obtain
\begin{align*}
\int \frac{3x+2}{x^2 (x-1) }  \d{x}&=\int \frac{-5x-2}{x^2} +\frac{5}{x-1}\d{x}\\
                                   &=\int -\frac{5}{x}-\frac{2}{x^2} +\frac{5}{x-1}\d{x}\\
                                   &=-5 \ln(x)+2\inverse{x}+5
                                     \ln(x-1)+C.
\end{align*}
\end{myexample}
Notice in this example that we really needed to look at the most
general possible polynomial of degree one less than the denominators
in the partial fractions decomposition; we would not have obtained the
correct answer otherwise.

\begin{embeddedproblem}{}
   Show that if we tried the following decomposition
$$
\frac{3x+2}{x^2 (x-1) }=\frac{A}{x^2} +\frac{B}{x-1}
$$
then no values for $A$ or $B$ would satisfy this.  [This says that we really
need to use the most general polynomial of degree one less than the
denominators in the partial fractions decomposition.]
\end{embeddedproblem}

Suppose we had 
$\int \frac{25}{(x-2)^2 (x^2+1)}  \d{x}$
Applying our partial fractions decomposition, we get
\begin{align*}
\frac{25}{(x-2)^2 (x^2+1) }&=\frac{Ax+B}{(x-2)^2} +\frac{Cx+D}{x^2+1}\\
                           &=\frac{(Ax+B)(x^2+1)+(Cx+D) (x-2)^2}{(x-2)^2 (x^2+1) }\\
                           &=\frac{(A+C) x^3+(B-4C-D)
                             x^2+(A+4C+4D)x+(B-4D)}{(x-2)^2 (x^2+1) }.
\end{align*}
This leads to 
$$
\hfill A+C=0,\hfill   B-4C-D=0,\hfill   A+4C+4D=0,\hfill      B-4D=25\hfill
$$
Solving this (any way you can), we get $
A=-4, B=13, C=4,$ and $ D=-3.$
Getting back to the integral, we have
\begin{align*}
\int \frac{25}{(x-2)^2 (x^2+1) }  \d{x}&=\int \frac{-4x+13)}{(x-2)^2}
                                         \d{x}+\int \frac{4x-3}{x^2+1}
                                         \d{x}\\ 
&=\int \frac{-4x+13}{(x-2)^2}  \d{x}+4\int \frac{x}{x^2+1} \d{x}
                                                    -3\int
                                                    \frac{1}{x^2+1}
                                                    \d{x}\\
&=\int \frac{-4x+13}{(x-2)^2}  \d{x}+2 \ln(x^2+1)-3\inverse\tan(x)+C
\end{align*}
Notice that the last wo integrals were pretty straightforward.  This
first integral is more problematic.  There is a little algebraic trick
which will allow us to break up that integral even further. 
\begin{align*}
\in\frac{t -4x+13}{(x-2)^2}  \d{x}&=\int
                                    \frac{-4((x-2)+2)+13}{(x-2)^2}
                                    \d{x} \\
                                  &=\int \frac{-4(x-2)-8+13}{(x-2)^2}
                                    \d{x} \\
                                  &=\int \frac{-4(x-2)+5}{(x-2)^2}
                                    \d{x} \\
                                  &=-4\int \frac{1}{x-2}  \d{x}+5\int
                                    \frac{1}{(x-2)^2}  \d{x} \\
                                  &=-4 \ln(x-2)-5(x-2)^{-1}
\end{align*}
Putting this all together, we finally have
$$
\int \frac{25}{(x-2)^2 (x^2+1)}  \d{x}=-4 \ln(x-2)-5\inverse{(x-2)}+2
\ln(x^2+1)-3 \inverse{\tan}(\theta)+C
$$

The trick that we employed above is something you might want to
remember.  It allows us to take our partial fractions decomposition a
bit further.  For example, if we start with 
$$
\int \frac{x^3}{(2x^2+1)^2}   \d{x}
$$
we could rewrite this as 
\begin{align*}
\int \frac{x^3}{(2x^2+1)^2}   \d{x}&=\int \frac{\frac{1}{2}
                                     x(2x^2)}{(2x^2+1)^2}   \d{x} \\
  &=\int \frac{\frac{1}{2} x(2x^2+1-1)}{(2x^2+1)^2}   \d{x} \\
  &=\frac{1}{2} \int \frac{x}{2x^2+1}  \d{x}-\frac{1}{2} \int
    \frac{x}{(2x^2+1)^2}   \d{x}. 
\end{align*}
Integrating this would be a matter of utilizing the substitution
$y=2x^2+1.$  Of course, this could have been done earlier to produce 
\begin{align*}
\int \frac{x^3}{(2x^2+1)^2}   \d{x}&=\int \frac{x^2}{(2x^2+1)^2}  x
                                     \d{x} \\
  &=\int \frac{\frac{1}{2} (y-1)}{y^2}  \left(\frac{1}{4} \d{y}\right) \\
  &=\frac{1}{8} \left[\int \frac{1}{y}  \d{y}-\int \frac{1}{y^2}   \d{y}\right]
\end{align*}
The point is that you have more tools at your disposal to transform
integrals you don't know how to do into ones that you recognize.

Before we send you off to practice such integrals, we want to present you
one more option.  You have probably noticed that there is more algebra
involved than calculus.  Some of it is unavoidable, but there are some
shortcuts.  For example, finding $A,B,C,$ and $D$ in the original partial
fractions decomposition required solving $4$ equations in $4$ unknowns.
Again, this was straightforward, but it turns out that there is a
shortcut for this similar to the HCUM.  We will present it here.  \textcolor{red}{Bud
- at this point, how much of the paper you wrote should we put in?
Will putting it in be too much at this point?  Should be some practice
exercises first?  Should it be appendicized?  I'm afraid I've already
written too much on partial fractions decompositions.   I need your
fresh view on this. }

\noindent{\bf{}Areas and Integration}\\
Recall, much earlier, that we examined the catenary (hanging chain).
We showed that the catenary was not a parabola as Galileo had
mentioned.  The idea of a catenary was (and is) important in
engineering in the construction of arches.  Masons before Galileo's
time were aware of the fact that the ``best'' arches were in the shape
of an inverted catenary, and they constructed wooden frames for their
stone utilizing a hanging chain for the correct shape.  This can be
seen in the following example.

\begin{myexample}
  Consider one half of a stone arch bridge drawn below.\\
  \centerline{\includegraphics*[height=2in,width=2in]{../Figures/StoneArch1}}

  We will draw the positive $y-$axis downward and will focus on the
  forces at a generic point $P$ with coordinates $(x,y).$  The problem is
  to find the curve $y=y(x)$ so that the vertical component of the
  tangential force $T$ at $P$ is equal to the weight of the bridge
  from $0$ 
  to $x.$  If we do this, then that means that the weight of the bridge
  will be directed toward the base of the bridge.  With this in mind,
  let $A(x)$ denote the area of the side section of the bridge from $0$ to
  $x.$  We will also let h denote the (constant) magnitude of the
  horizontal force along the length of the bridge and $w$ be the weight
  density of the stone (per cross sectional area).\\
  \centerline{\includegraphics*[height=2in,width=2in]{../Figures/StoneArch2}}
   With all of this set up, what we really want is the horizontal
   component of the force $T$ to be $h$ and the vertical component to be
   the weight of the bridge from $0$ to $x,$ namely, $wA.$  This leads to
   the following picture. \\
   \centerline{\includegraphics*[height=2in,width=2in]{../Figures/StoneArch3}}
   This leads to the slope of the tangent line at $P$ equaling
   $\frac{wA}{h},$ so we get the differential equation 
$$
\dfdx{y}{x}=\frac{w}{h} A.
$$
The problem here is that we don't know what $A$ is.  However, a
moment's thought tells us that we know what $\d{A}$ is.  If we
increase $x$ to $x+\d{x},$ then we can see that $\d{A}=y\d{x}.$\\
\centerline{\includegraphics*[height=2in,width=2in]{../Figures/StoneArch4}}

Thus, if we differentiate the above equation, we get that the arch
should satisfy the differential equation 
$$
\dfdxn{y}{x}{2}=\frac{w}{h}  \dfdx{A}{x}=\frac{w}{h}  y.
$$
\end{myexample}

\begin{embeddedproblem}{}
\begin{enumerate}[label={\bf{}(\alph*)}]
\item 	Show that $y=Ae^{\sqrt{\frac{w}{h}} x}+Be^{-\sqrt{\frac{w}{h}} x}$
satisfies the above differential equation.
\item Show that if $y(0)=\sqrt{\frac{h}{w}}$  and  $\left.\dfdx{y}{x}\right|_{x=0}=0,$ then $A=B=\frac{1}{2}\sqrt{\frac{h}{w}}.$  Compare this to the equation of the catenary in Problem in Context\footnote{\textcolor{red}{PIC number is hard-coded; needs to be a \\ref}} \# 53.
\end{enumerate}
\end{embeddedproblem}

While it is interesting that the above example illustrates the
importance of catenaries in engineering, the important thing for us is
that $\d{A}=y\d{x}.$   Thus  
$$
\int y \d{x}=\int \d{A}=A+C
$$
which relates integrals to areas.  This relationship was known to both
Newton and Leibniz, as well as some of their predecessors.  It is of
such importance that it is known as the {\bf{}Fundamental Theorem of
  Calculus. }

To make this a bit more concrete, let's set this up a bit more
carefully.  Suppose we start with a curve $y=y(x)$ where $y\ge0$ for
$a\le x\le b,$ and we wish to find the area of the region bounded by
the curves $y=y(x), y=0, x=a,$ and $x=b.$ \\
\centerline{\includegraphics*[height=2in,width=2in]{../Figures/FTCalc1}}

We used the notation $\int y\d{x}$ to denote a sum of differentials
and called this integration with the idea that it was the opposite of
differentiation and provided a notation for the antiderivative to the
function $y.$  In examining a situation like we have with areas above,
we will modify our notation a bit to
$$
\int_{x=a}^b y \d{x}
$$
to emphasize that we are adding differentials from $x=a$ to $x=b,$ so that
$\int_{x=a}^by \d{x}$ represents a specific number (in this case area),
whereas $\int y \d{x}$ represents an antiderivative which is a family of
functions differing by arbitrary constants.  With this in mind,
$\int_{x=a}^by \d{x}$ is called the definite integral of $y \d{x}$ from $x=a$ to
$x=b$ and $\int y \d{x}$ is called the indefinite integral of $y \d{x}.$  So, the
definite integral is a specific number while the indefinite integral
is a family of functions.  If using almost the same notation seems
confusing at this point, it is brought together by the Fundamental
Theorem of Calculus.
\begin{mytheorem}[\bf{}The Fundamental Theorem of Calculus]
Suppose that $\dfdx{A}{x}=y$ (so that $A$ is an antiderivative of $y$).  Then
$$
\int_{x=a}^b y \d{x}=A(b)-A(a).
$$
\end{mytheorem}

As was mentioned before, this result was known to both Newton and
Leibniz and was first published in a short paper by Leibniz in $1693.$
The idea is so surprisingly simple and profound, that it warrants an
(apocryphal) story.  When Leibniz was a young diplomat in Paris, he
sought out Christian Huygens who, at the time, was one of the most
brilliant mathematicians in the world.  Leibniz credited much of his
mathematical growth in Paris to Huygens.  The idea for summing
differences came to Leibniz from a problem given to him by Huygens in
$1672,$ namely to compute the following sum
$$
\frac{1}{1\cdot 2}+\frac{1}{2\cdot 3}+\frac{1}{3\cdot 4} + \ldots{}
$$
Leibniz recognized that this could be rewritten as
$$
\left(\frac{1}{1}-\frac{1}{2}\right)+\left(\frac{1}{2}-\frac{1}{3}\right)+\left(\frac{1}{3}-\frac{1}{4}\right)+\ldots{}
$$
and noted that nearly all of the fractions cancelled out leaving
$1-\frac{1}{\infty{}}=1-0=1.$

The importance of this was not the result, but the technique, namely
that the sum of differences is equal to the difference of the
extremes.  Applying this same principle to the Fundamental Theorem of
calculus, we have
\begin{align*}
\dfdx{A}{x}&=y \\
  y \d{x}&=\d{A} \\
  \int_{x=a}^by \d{x}&=\int_{x=a}^b\d{A}=A(b)-A(a).
\end{align*}

This was so quick, that it probably deserves a picture.  Consider the
picture below of the two functions $y$ and $A$ related by the fact
that
$\dfdx{A}{x}=y.$\\
\centerline{\includegraphics*[height=2in,width=2in]{../Figures/FTCalc2}}
Notice that since $y \d{x}=\d{A},$ then this says that the area of the
box with width $\d{x}$ and height $y$ in the first diagram is the same
as the length $\d{A}$ of the segment in the second diagram.  Od
course, these represent corresponding generic boxes and segments.  If
we add all of these together, the sum of the areas of the boxes in the
first diagram will provide the area of the region; the sum of the
lengths of the segments in the second diagram will provide the length
of the segment from $A(a)$ to $A(b).$ In other words, we have
$$
\int_{x=a}^b y \d{x}=\int_{x=a}^b\d{A}=A(b)-A(a)
$$ 
which is what the Fundamental Theorem of Calculus says.

It is hard to overestimate the importance of this result. (Not
everything is called a Fundamental Theorem.)  In many applications, we
will have a need to add all of the differentials which may or may not
represent areas.  After we set up this sum, computation will bring us
back to our old problem of computing antiderivatives.

%%% Local Variables:
%%% mode: latex
%%% TeX-master: "IntCalc"
%%% End:
